\uuid{3OCj}
\exo7id{754}
\titre{exo7 754}
\auteur{gourio}
\organisation{exo7}
\datecreate{2001-09-01}
\isIndication{false}
\isCorrection{false}
\chapitre{Fonctions circulaires et hyperboliques inverses}
\sousChapitre{Fonctions circulaires inverses}
\cpge{mpsi-mp2i.calcul-algebrique-trigonometrie.a, pcsi.calcul-algebrique-trigonometrie.a, ptsi.calcul-algebrique-trigonometrie.a, mpsi-mp2i.nombres-reels-suites.d, pcsi.nombres-reels-suites.c, ptsi.nombres-reels-suites.c, mpsi-mp2i.analyse-asymptotique.c, pcsi.analyse-asymptotique.c, ptsi.analyse-asymptotique.c, bcpst1.methodes-calcul, bcpst1.suites-reelles}
\module{Analyse}
\niveau{L1}
\difficulte{}

\contenu{
\question{
\'Etudier la suite $(u_{n})_{n\in \Nn}$ d\'{e}finie par :
$$\forall n\in \Nn,u_{n}=\sum\limits_{k=1}^{n}\arctan \frac{1}{k^{2}-k+1}.$$
On montrera qu'elle converge (vers $\ell$) et on \'{e}valuera $\lim_{n\rightarrow \infty }n(u_{n}-\ell)$.

\emph{Indication} : que vaut $\arctan a-\arctan b$ ?
}
}
